% Abstract - guideline 4.1.6: optional, 120-150 words recommended,
% describes the problem and results, with key words.

\chapter*{Abstract}
\addcontentsline{toc}{chapter}{Abstract}

% Word count checked against guideline 4.1.6 (120-150 words recommended);
% re-check if this text is edited further.

Autonomous mobile robots are typically developed in simulation before
deployment on hardware, yet the transfer to a real platform routinely
exposes gaps that simulated testing alone does not predict. This
thesis documents the migration of a ROS~2 Nav2 navigation stack from
Gazebo simulation onto QCar, Quanser's Ackermann-steered vehicle.
Limited onboard storage, together with a ROS~2 Dashing/Humble
incompatibility on the robot's onboard computer, motivated a custom
TCP/JSON bridge. The thesis presents the bridge's design, the
calibration of an open-loop throttle model against real actuator
behaviour, the diagnosis of hardware defects and building an
end-to-end full-stack navigation system, incorporating SLAM via
Cartographer, AMCL localisation and MPPI-based control. Validation
follows through closed-loop Nav2 goal-reaching tests on the vehicle,
with the observed sim-to-real discrepancies quantified.

\vspace{0.5cm}
\noindent\textbf{Key words:} ROS~2, Nav2, SLAM, Cartographer, AMCL,
Ackermann steering, Model Predictive Path Integral control,
sim-to-real transfer, Quanser QCar
